\section{Conclusion}
Adaptive interaction graphs offer a direct way to spend computation where a particle simulator needs it. The central difficulty is deciding what that need means. A residual head predicts error under the current model; an allocator must predict whether adding particular interactions will reduce that error. Our cached controller and explicit pair budget make this distinction testable while separating the benefit of graph exposure from the benefit of edge placement.

Training on expanded graphs improves fixed-random rollouts in every paired Goop and WaterDrop seed, but residual allocation remains conditional. WaterDrop's observed-test risk ordering reverses after exposure without a consistent autonomous interaction. Sand sharpens the distinction: exposure narrows the observed risk deficit while worsening cached-risk rollout error and its gap to random in every seed; native base also beats random expansion in every seed of both arms. Useful graph adaptation therefore requires evidence at three levels: the model can use the added messages, the score predicts their intervention benefit, and the resulting feedback remains reliable. Residual accuracy and smaller realized rollout graphs alone do not establish those properties.
